% 本文件由 LaTeX 排版 Agent 根据有效 translation.json 自动生成。
% author=Polycarp; work_id=0136; title=波利卡普致斐理伯人书

\chapter{波利卡普致斐理伯人书}

% structure: p0001 source=h1 target=omitted reason=page-title-already-rendered-by-parent

% structure: p0002 source=h2 target=section reason=relative-heading-rank; base=h2

\section{问候}

\Person{波利卡普}，以及与他一起的长老们，致书给在\Place{斐理伯}寄居的天主教会：愿来自全能天主并我们的救主主耶稣基督的怜悯与平安，多多加给你们。\par

% structure: p0004 source=h2 target=section reason=relative-heading-rank; base=h2

\section{第一章 称赞斐理伯人}

我因你们在我们的主耶稣基督内大大欢喜，因为你们效法了真实的爱（如同天主所显示的），并且相称地陪伴了那些被锁链捆锁的人，这锁链是圣人相称的装饰，实为天主和我们的主之真选民的冠冕；并且因为你们信德的坚固根基，正如从前在那些日子里所传讲的\BibleRef{斐 1:5}，至今仍存，并向我们的主耶稣基督结果实，他为了我们的罪甚至受苦至死，\Quote{但是天主却解除了坟墓的束缚，使他从死人中复活。} \BibleQuote{你们虽然没有见过他，却信了他，因信而有说不出的、充满光荣的喜乐；}\BibleRef{伯前 1:8}许多人都渴望进入那喜乐，知道\BibleQuote{你们是因恩宠得救，不是出于行为，}\BibleRef{弗 2:8}而是藉着耶稣基督，靠着天主的旨意。\par

% structure: p0006 source=h2 target=section reason=relative-heading-rank; base=h2

\section{第二章 劝勉德行}

\BibleQuote{所以，你们要束上腰，}\BibleRef{伯前 1:13} \BibleQuote{怀着敬畏之心侍奉主，} 并以真理侍奉，如同那些抛弃了众百姓虚妄空洞的言谈和谬误的人，并且\BibleQuote{信了那使我们的主耶稣基督从死人中复活、并赐给他荣耀的那一位，}\BibleRef{伯前 1:21}及在他右边的宝座。天上地下的万有\BibleRef{伯前 3:22}都服属于他。众灵都事奉他。他要来审判活人死人。\BibleRef{宗 17:31}凡不信他的人，天主必追讨他的血。但那使他从死人中复活的那一位，也必使我们复活，如果我们遵行他的旨意，行走在他的诫命中，并爱他所爱的，保守自己远离一切不义、贪婪、爱财、恶言、假见证；\BibleQuote{不以恶报恶，不以辱骂还辱骂，}\BibleRef{伯前 3:9}不以打还打，不以咒骂还咒骂，却要记念主在他的教导中所说的：\Quote{你们不要判断，免得你们被判断；\BibleRef{玛 7:1}你们宽恕，就必蒙宽恕；你们要慈悲，好蒙慈悲；\BibleRef{路 6:36}你们用什么量器量给人，也必用什么量器量给你们；\BibleRef{玛 7:2}} 并且又说：\BibleQuote{贫穷的人有福了，为义受迫害的人有福了，因为天国是他们的。}\par

% structure: p0008 source=h2 target=section reason=relative-heading-rank; base=h2

\section{第三章 表达个人的不配}

弟兄们，我写这些关于正义的事给你们，并非因为我自取什么，而是因为你们邀请了我。因为无论是我，还是其他任何这样的人，都不能及得上受颂扬的、满被光荣的保禄的智慧\BibleRef{伯后 3:15}。他在你们中间时，对那些当时活着的人，准确而坚定地传讲了真理之道。当他离开你们时，又写了一封信给你们。你们若仔细研读，便会发现它能使你们在所领受的信德上得到建立；这信德后有望德相随，前有对天主、对基督及对邻人的爱德为前导，\BibleQuote{它是我们众人的母亲。}\BibleRef{迦 4:26}因为任何人若内里拥有这些恩宠，便满全了正义的诫命，凡有爱德的人远离一切罪恶。\par

% structure: p0010 source=h2 target=section reason=relative-heading-rank; base=h2

\section{第四章 各种劝勉}

\BibleQuote{贪爱钱财是万恶的根源。}\BibleRef{弟前 6:10}因此，既然知道\BibleQuote{我们什么都没有带到世界上来，也不能带走什么，}\BibleRef{弟前 6:7}就让我们以正义的武器武装自己；\BibleRef{弗 6:11}并首先教导我们自己行走在主耶稣基督的诫命中。其次，教导你们的妻子在赐给她们的信德、爱德和纯洁中行走，以一切真理温柔地爱自己的丈夫，并以一切贞洁同等爱所有人；还要教导她们的儿女在认识天主和敬畏天主中成长。教导寡妇们在有关主的信德上要明智，不断祈祷\BibleRef{得前 5:17}为众人，远离一切毁谤、恶言、假见证、贪财及各种邪恶；知道她们是天主的祭台，他清楚洞察一切，没有什么能向他隐藏——无论是心思、意念还是心中任何隐秘之事。\par

% structure: p0012 source=h2 target=section reason=relative-heading-rank; base=h2

\section{第五章 执事、青年及童贞女的职责}

\BibleQuote{天主是轻慢不得的，}\BibleRef{迦 6:7}我们应当以相称于他的诫命和荣耀的方式行走。同样，执事们应当在他的正义面前无可指摘，作为天主和基督的仆人，而非人的仆人。他们不可诽谤、一口两舌\BibleRef{弟前 3:8}、贪爱钱财，而要在一切事上节制、慈悲、勤勉，按主的真理行走，主曾是众人的仆人\BibleRef{玛 20:28}。如果我们在这现世中讨他喜悦，我们也将获得那将来的世界，正如他应许我们的：他将使我们从死人中复活，并且如果我们活在他面前，\BibleQuote{我们也要与他一同为王，}\BibleRef{弟后 2:12}只要我们相信。同样，青年男子也应当在一切事上无可指摘，特别谨慎保持纯洁，如同用缰绳约束自己，远离各种邪恶。因为将他们与世上的情欲割断是好的，因为\BibleQuote{情欲与灵魂交战；}\BibleRef{伯前 2:11}并且\BibleQuote{无论是行淫的、柔弱的、与男子同寝的，都不能承受天主的国，}\BibleRef{格前 6:9}那些行为不合宜、不体面的人也是如此。因此，必须戒绝这一切，顺从长老和执事，如同顺从天主和基督。童贞女也当以无可指摘的纯洁良心行走。\par

% structure: p0014 source=h2 target=section reason=relative-heading-rank; base=h2

\section{第六章 长老及其他人的职责}

并且让司铎们对众人有同情心和怜悯，领回那些迷失的人，探望所有病人，不忽略寡妇、孤儿或穷人，但常要\Quote{准备在神和人面前做合适的事；}\BibleRef{罗 12:17}戒除一切愤怒、偏袒和不义的判断；远离一切贪财，不轻易相信对任何人的[恶言]控告，在判断上不严厉，因为我们都知道自己欠着罪债。如果我们求主宽恕我们，我们也应当宽恕别人；\BibleRef{玛 6:12}因为我们是在主和天主的眼前，并且\Quote{我们众人必须出现在基督的审判台前，每人要交出自己的账。}\BibleRef{罗 14:10}因此，我们要在恐惧中事奉祂，并怀着一切尊敬，正如祂亲自命令我们的，也如那传福音给我们的宗徒和预言主来临的先知[都同样教导我们]。我们要热切追求善事，避免犯罪的机会、假弟兄和那些假冒主名、将虚妄之人引入错误的人。\par

% structure: p0016 source=h2 target=section reason=relative-heading-rank; base=h2

\section{第七章　避免幻象论者，并持守斋戒和祈祷}

\Quote{因为凡不承认耶稣基督降生成人的，就是假基督；}\BibleRef{若一 4:3}凡不承认十字架见证的，是出于魔鬼；凡随自己的私欲曲解主的圣言，并说没有复活和审判的，他是撒殚的长子。所以，我们要抛弃许多人的虚妄和假道理，返回那从起初\BibleRef{犹 1:3}传给我们的圣言；\Quote{警醒祈祷，}\BibleRef{伯前 4:7}并持守斋戒；在恳求中向那洞察万有的天主祈求\Quote{不要让我们陷入诱惑，}\BibleRef{玛 6:13}正如主所说：\Quote{心神固然愿意，但肉体却软弱。}\BibleRef{玛 26:41}\par

% structure: p0018 source=h2 target=section reason=relative-heading-rank; base=h2

\section{第八章　持守望德和忍耐}

因此，我们要不断持守我们的望德和我们正义的凭据，就是耶稣基督，\Quote{祂亲身在木头上承担了我们的罪过，}\BibleRef{伯前 2:24}\Quote{祂没有犯过罪，口中也找不到诡诈，}\BibleRef{伯前 2:22}但为了我们忍受了一切，好使我们因祂而生活。\BibleRef{若一 4:9}因此，我们要效法祂的忍耐；如果我们为祂的名受苦，\BibleRef{宗 5:41}就让我们光荣祂。因为祂亲自给我们立了榜样，\BibleRef{伯前 2:21}我们也相信事实如此。\par

% structure: p0020 source=h2 target=section reason=relative-heading-rank; base=h2

\section{第九章　劝勉忍耐}

因此，我劝你们大家，要顺从正义的言语，并要操练一切忍耐，就如你们曾[亲眼]看见的，不仅在蒙福的\Person{依纳爵}、\Person{若西默}和\Person{鲁富}身上，而且在你们中间的其他人身上，以及在\Person{保禄}本人和其余宗徒身上。[你们要这样做]确知所有这些人都没有白跑\BibleRef{斐 2:16}，而是怀着信德和义德，并且他们[如今]在主面前各得其所，与祂一同受苦。因为他们不爱这现世，却爱那为我们死了，并为我们从死者中复活的祂。\par

% structure: p0022 source=h2 target=section reason=relative-heading-rank; base=h2

\section{第十章　劝修德行}

因此，你们要站稳在这些事上，并效法主的榜样，在信德上坚定不变，友爱弟兄\BibleRef{伯前 2:17}，彼此依恋，在真理中结合，彼此交往中表现出主的温和，不轻视任何人。你们若可行善，不要迟延，因为\Quote{施舍能救人脱离死亡。}\BibleRef{多 4:10}你们众人要彼此顺服\BibleRef{伯前 5:5}，\Quote{在外邦人中品行端正，}\BibleRef{伯前 2:12}好使你们因善工得称赞，并且主的名不致因你们受亵渎。但藉谁主的名受亵渎，那人就有祸了！\BibleRef{依 52:5}因此，要教众人节制，并在你们自己的品行上也显示出来。\par

% structure: p0024 source=h2 target=section reason=relative-heading-rank; base=h2

\section{第十一章　为瓦冷斯表达忧伤}

我为\Person{瓦冷斯}非常忧伤，他从前是你们中的一位司铎，因为他如此不明白赐予他的[在教会内的]地位。因此，我劝你们戒除贪财，并要贞洁和诚实。\Quote{要戒避一切邪恶的形态。}\BibleRef{得前 5:22}如果有人连在这些事上都不能自制，他怎能命令别人呢？如果有人不戒避贪财，他必被拜偶像所玷污，并被当作外邦人来审判。但我们中谁不知道主的审判呢？\Quote{我们不知道圣人将要审判世界吗？}\BibleRef{格前 6:2}正如\Person{保禄}所教导的。但我从未在你们中间看到或听到这样的事，在你们中间蒙福的\Person{保禄}曾劳苦，你们在他的书信开头受到称赞。因为他曾在所有那些当时唯独认识主的教会中夸耀你们；但我们[斯米纳人]尚未认识祂。因此，弟兄们，我为他(\Person{瓦冷斯})和他的妻子深感忧伤；愿主赐予他们真正的悔改！对于这件事，你们要适度，并且\Quote{不要以他们为仇敌，}\BibleRef{得后 3:15}而要挽回他们，如同挽回受苦和迷路的肢体，好使你们拯救整个身体。因为这样行，你们将彼此建立。\BibleRef{格前 12:26}\par

% structure: p0026 source=h2 target=section reason=relative-heading-rank; base=h2

\section{第十二章　劝勉各种恩惠}

因为我确信你们精通圣经，没有任何事向你们隐藏；但这\OriginalTerm{特权}{privilege}尚未赐给我。在这些圣经中如此说：\Quote{你们生气，却不要犯罪；}又说：\Quote{不可让太阳在你们的愤怒中落下。}\BibleRef{弗 4:26}那记住这话的人是有福的，我相信你们正是如此。但愿我们的主耶稣基督的天主和父，以及耶稣基督自己——祂是天主子，我们永远的\OriginalTerm{大司祭}{High Priest}——在信德和真理、以及一切温良、柔和、忍耐、宽容、坚忍和纯洁中建立你们；并愿祂在祂的圣人们中赐给你们一份产业，也赐给我们和你们，以及天下所有将要相信我们的主耶稣基督和祂的父的人，那位\Quote{使祂从死者中复活}的父。\BibleRef{迦 1:1}请为所有圣人祈祷。也为君王、\BibleRef{弟前 2:2}掌权者、首领，以及那些迫害和仇恨你们的人祈祷，\BibleRef{玛 5:44}并为十字架的仇敌祈祷，使你们的果实向众人显明，并使你们在祂内得以完全。\par

% structure: p0028 source=h2 target=section reason=relative-heading-rank; base=h2

\section{第十三章　关于信件的传递}

你们和\Person{依纳爵}都写信给我，说如果有人[从这里]去叙利亚，应随身携带你们的信；我若找到合适的机会，无论是亲自还是通过他人代我行事，都会遵行，以满足你们的愿望。\Person{依纳爵}写给我们的信函，以及我们手中所有的其他[信函]，都已按你们的要求寄给你们。这些信函附在本信之后，你们可以从中大得益处；因为它们论及信德和忍耐，以及一切有助于在主内建立德行的事。关于\Person{依纳爵}本人以及与他同在的人，你们若得到任何更确切的消息，请惠告我们。\par

% structure: p0030 source=h2 target=section reason=relative-heading-rank; base=h2

\section{第十四章　结语}

这些事我藉\Person{克雷森斯}写信给你们，我至今向你们推荐他，现在仍推荐。因为他在我们中间行为无可指责，我相信在你们中间也是如此。此外，当他的姊妹来到你们那里时，请你们尊重她。愿你们在主耶稣基督内平安。愿恩宠与你们众人同在。阿们。\par

\SourceRecord{Translated by Alexander Roberts and James Donaldson.；Ante-Nicene Fathers；Vol. 1.；Edited by Alexander Roberts, James Donaldson, and A. Cleveland Coxe.；Buffalo, NY: Christian Literature Publishing Co.,；1885.；<http://www.newadvent.org/fathers/0136.htm>.}
